% FORMALIZACION_REUNIDA. Desarrollo ampliado de las tres composiciones de
% EXPLORACION_AMBIVALENCIA_ESPEJO_DIAGONAL_20260918, con dependencias de X.
% Las coordenadas pi, phi, e, alpha y el registro K son salidas generadas.
% Este capítulo define lectores posteriores y no altera el generador TPK.
\chapter{Holonomía nonádica, reciprocidad y recuperación de lectores geométricos}
\label{hrr:capitulo}

La conexión nonádica, la incidencia areal--volumétrica y la reciprocidad
radial intervienen en realizaciones diferentes del mismo estado enriquecido.
La conexión conserva la sucesión orientada de las nueve fases y prolonga
la memoria al retornar. La incidencia registra las transiciones entre dos
modos de lectura. La reciprocidad radial intercambia las dos orientaciones
de una forma geométrica de área fijada. Su composición se expresa mediante
tres resultados precisos: la invariancia de un cociente de marcas a toda
profundidad, la factorización de una deformación de determinante uno y la
reconstrucción de la anisotropía a partir de una longitud normalizada y su
orientación.

El orden de construcción permanece
\[
 \begin{gathered}
 \mathrm{APP}\longrightarrow\mathrm{TRIT}\longrightarrow\mathrm{TPK}
 \longrightarrow\text{estado enriquecido},\\
 \text{estado enriquecido}
 \longrightarrow\text{estructura discreta conjunta del continuo},\\
 \text{estructura discreta conjunta del continuo}
 \longrightarrow\text{lectores aquí compuestos}.
 \end{gathered}
\]
Las coordenadas \(\pi,\varphi,e,\alpha\) y el registro \(K\) que aparecen
en las fórmulas son las salidas de la generación regional precedente. El
calendario trítico ya determina los pares consecutivos de los modos
aditivo y multiplicativo. El lenguaje de matrices, funciones hiperbólicas
y menores exteriores explicita a continuación las realizaciones de esa
estructura generada.

\section{Incidencia modal y marcas regionales sobre retornos finitos}
\label{hrr:incidencia}

\subsection{Enumeración de los retornos entre hojas}

El bloque temporal \((+,-,0)\) induce la sucesión modal
\((\Sigma,\Pi,\Pi)\). Sus pares consecutivos, incluido el enlace de retorno,
son \((\Sigma,\Pi)\), \((\Pi,\Pi)\) y \((\Pi,\Sigma)\). La incidencia
primitiva, en el orden areal--volumétrico, es por tanto
\begin{equation}
 F_{\mathrm{av}}=
 \begin{pmatrix}0&1\\1&1\end{pmatrix}.
 \label{hrr:incidencia-matriz}
\end{equation}
Esta matriz registra los pares del calendario. El envolvente simbólico
de memoria uno que determina conserva esa incidencia; las restricciones
adicionales de fase, orientación y registro permanecen en el estado TPK.
Así se distinguen el recuento del envolvente y la supervivencia de una
historia enriquecida completa.

Sean \(f_0=0\), \(f_1=1\) y
\(f_{n+1}=f_n+f_{n-1}\) para \(n\geq1\). Estos enteros expresan aquí
los recuentos producidos por las potencias de la incidencia. La
identificación convencional de esta sucesión conserva su carácter
posterior a la construcción del calendario.

\begin{proposition}[Potencias de la incidencia y retornos diagonales]
\label{hrr:potencias}
Para todo \(n\geq1\),
\begin{equation}
 F_{\mathrm{av}}^n=
 \begin{pmatrix}f_{n-1}&f_n\\f_n&f_{n+1}\end{pmatrix}.
 \label{hrr:potencias-formula}
\end{equation}
En particular, los retornos de longitud \(n\) a las hojas areal y
volumétrica tienen multiplicidades \(f_{n-1}\) y \(f_{n+1}\),
respectivamente.
\end{proposition}
\begin{proof}
Para \(n=1\), la fórmula es la definición de
\(F_{\mathrm{av}}\), pues \(f_2=1\). Supuesta válida para \(n\),
la multiplicación por la derecha proporciona
\[
 \begin{pmatrix}f_{n-1}&f_n\\f_n&f_{n+1}\end{pmatrix}
 \begin{pmatrix}0&1\\1&1\end{pmatrix}
 =\begin{pmatrix}f_n&f_{n-1}+f_n\\
 f_{n+1}&f_n+f_{n+1}\end{pmatrix}
 =\begin{pmatrix}f_n&f_{n+1}\\f_{n+1}&f_{n+2}\end{pmatrix}.
\]
La inducción prueba la identidad. Para una matriz de incidencia con
entradas cero o uno, la regla de multiplicación suma las sucesiones de
transiciones compatibles entre dos posiciones. Por inducción sobre la
longitud, la entrada \((i,j)\) de la potencia cuenta las sucesiones de
longitud \(n\) de \(i\) a \(j\). Las dos entradas diagonales cuentan
los retornos declarados.
\end{proof}

\subsection{Marcación regional del soporte de retorno}

Sean \(e_{\mathrm{ar}}=(1,0)^{\mathsf T}\) y
\(e_{\mathrm{vol}}=(0,1)^{\mathsf T}\). Para una coordenada positiva
\(c\), producida antes de esta evaluación, definimos el lector de marca
\[
 D_c=c\,e_{\mathrm{ar}}e_{\mathrm{ar}}^{\mathsf T}
       +e_{\mathrm{vol}}e_{\mathrm{vol}}^{\mathsf T}
     =\operatorname{diag}(c,1).
\]
La diagonalidad conserva las dos hojas; la entrada volumétrica expresa
su normalización unitaria. La marca areal \(c\) se aplica una sola vez,
después del recuento de retornos. Su evaluación es
\begin{equation}
 \Theta_{c,n}=
 \frac{\langle e_{\mathrm{ar}},D_c F_{\mathrm{av}}^n
 e_{\mathrm{ar}}\rangle}
 {\langle e_{\mathrm{vol}},F_{\mathrm{av}}^n
 e_{\mathrm{vol}}\rangle}
 =c\,\frac{f_{n-1}}{f_{n+1}},\qquad n\geq1.
 \label{hrr:marca}
\end{equation}
La elección \(c=\pi\) coincide con el lector de frontera del retorno
areal--volumétrico. La elección \(c=e\) define el lector posterior
complementario considerado en este capítulo. En ambos casos se utiliza
el mismo soporte de retornos y la coordenada regional ya generada.

\begin{theorem}[Cociente de marcas independiente de la profundidad]
\label{hrr:cociente}
Para cualesquiera marcas positivas \(c,d\) y todo \(n\geq2\),
\begin{equation}
 \frac{\Theta_{c,n}}{\Theta_{d,n}}=\frac cd.
 \label{hrr:cociente-exacto}
\end{equation}
En particular,
\begin{equation}
 \frac{\Theta_{\pi,n}}{\Theta_{e,n}}=\frac\pi e,
 \qquad
 \lim_{n\to\infty}\Theta_{\pi,n}=\frac\pi{\varphi^2},
 \qquad
 \lim_{n\to\infty}\Theta_{e,n}=\frac e{\varphi^2}.
 \label{hrr:pi-e-phi}
\end{equation}
\end{theorem}
\begin{proof}
La recurrencia implica \(f_j>0\) para \(j\geq1\). Para \(n\geq2\),
el cociente \(r_n=f_{n-1}/f_{n+1}\) es positivo. Por
\eqref{hrr:marca}, las dos evaluaciones son \(cr_n\) y \(dr_n\),
de modo que su división cancela exactamente el mismo factor de retorno.
La identificación precedente de la coordenada regional de autoescala
con la razón estable de esta incidencia da
\(r_n\to\varphi^{-2}\). Al multiplicar por las marcas constantes
\(\pi\) y \(e\), se obtienen ambos límites. El cociente finito,
en cambio, ya tiene su valor exacto antes del paso al límite.
\end{proof}

La primera longitud tiene \(f_0=0\): ambas evaluaciones se anulan y su
cociente queda fuera del dominio de \eqref{hrr:cociente-exacto}. La
restricción \(n\geq2\) expresa una propiedad del recuento y evita
extender una simplificación a un denominador nulo.

\begin{center}
\begin{tabular}{c|ccccc}
\(n\)&2&3&4&5&6\\\hline
\(r_n\)&\(1/2\)&\(1/3\)&\(2/5\)&\(3/8\)&\(5/13\)\\
\(\Theta_{\pi,n}\)&\(\pi/2\)&\(\pi/3\)&\(2\pi/5\)&\(3\pi/8\)&\(5\pi/13\)\\
\(\Theta_{e,n}\)&\(e/2\)&\(e/3\)&\(2e/5\)&\(3e/8\)&\(5e/13\)
\end{tabular}
\end{center}
La tabla exhibe dos hechos simultáneos: cada lector cambia con la
profundidad, mientras que la comparación entre ambos permanece fija.
La invariancia corresponde a una relación entre lectores y no a la
inmovilidad del estado que sostiene la lectura.

\subsection{Cota exacta de convergencia a la razón estable}

La correspondencia ya establecida de la coordenada generada
\(\varphi\) con el autovalor positivo de
\(F_{\mathrm{av}}\) proporciona \(\varphi^2=\varphi+1\) y
\(\varphi>1\). Es posible cuantificar el límite sin introducir
una expansión decimal objetivo.

\begin{proposition}[Error del cociente de retornos]
\label{hrr:error-retorno}
Puesto \(q=-\varphi^{-2}\), para \(n\geq1\) se cumple
\begin{align}
 r_n&=\varphi^{-2}\frac{1-q^{n-1}}{1-q^{n+1}},
 \label{hrr:rn-exacto}\\
 r_n-\varphi^{-2}
 &=\varphi^{-2}
 \frac{q^{n-1}(q^2-1)}{1-q^{n+1}}.
 \label{hrr:rn-error}
\end{align}
En particular,
\begin{equation}
 |r_n-\varphi^{-2}|
 \leq\varphi^{-2}
 \frac{|q|^{n-1}(1-q^2)}{1-|q|^{n+1}}.
 \label{hrr:rn-cota}
\end{equation}
\end{proposition}
\begin{proof}
Sea \(\psi=-\varphi^{-1}\). De \(\varphi^2=\varphi+1\)
se deduce \(\psi^2=\psi+1\). La sucesión
\(b_n=(\varphi^n-\psi^n)/(\varphi-\psi)\)
satisface \(b_0=0\), \(b_1=1\) y la misma recurrencia de
\(f_n\). La inducción da \(b_n=f_n\). Al dividir la expresión
de \(f_{n-1}\) por la de \(f_{n+1}\), y usar
\(\psi/\varphi=q\), resulta \eqref{hrr:rn-exacto}.
La resta del límite produce \eqref{hrr:rn-error}. Finalmente,
\(|q|<1\), \(1-q^2>0\) y
\(|1-q^{n+1}|\geq1-|q|^{n+1}>0\) dan la cota.
\end{proof}

Esta fórmula constituye una lectura exacta de la velocidad de convergencia
de la incidencia ya generada. Las marcas reciben el error multiplicado
por \(c\), pero su cociente conserva la cancelación exacta del
teorema~\ref{hrr:cociente}.

\section{Composición multiplicativa del flujo de acoplamiento}
\label{hrr:flujo}

\subsection{Dirección generada y eliminación de la componente uniforme}

El registro dodecafásico \(K\), con su marcación, produce la dirección
\[
 u_K=P_3P_{11}K,\qquad u_K\neq0,\qquad
 \langle u_K,\mathbf1\rangle=0.
\]
Los proyectores son los operadores ya construidos en la realización
correspondiente del registro. En esta sección se fija su resultado y
se escribe
\[
 v=\frac{u_K}{\|u_K\|},\qquad
 H_K=\operatorname{diag}(v_1,\ldots,v_{12}).
\]
Entonces \(\sum_i v_i=0\), \(\sum_i v_i^2=1\) y
\(\operatorname{tr}H_K=0\). Para una marca positiva \(r\) y un
parámetro real \(t\), definimos
\begin{equation}
 g_{K,r}(t)=\exp\bigl(t\log(r)H_K\bigr)
 =\operatorname{diag}\bigl(r^{tv_1},\ldots,r^{tv_{12}}\bigr).
 \label{hrr:g-definicion}
\end{equation}
La dirección procede del registro; la marca determina la escala
logarítmica de la deformación. El parámetro \(t\) parametriza este
flujo geométrico. Una identificación con un reloj físico requiere la
realización temporal correspondiente.

\begin{theorem}[Homomorfismo de marcas y conservación del determinante]
\label{hrr:homomorfismo}
Para \(a,b>0\) y \(s,t\in\mathbb R\),
\begin{align}
 g_{K,ab}(t)&=g_{K,a}(t)g_{K,b}(t),
 \label{hrr:producto}\\
 g_{K,a/b}(t)&=g_{K,a}(t)g_{K,b}(t)^{-1},
 \label{hrr:cociente-flujo}\\
 g_{K,r}(s+t)&=g_{K,r}(s)g_{K,r}(t),
 \label{hrr:grupo-t}\\
 \det g_{K,r}(t)&=1.
 \label{hrr:determinante}
\end{align}
Además, \(g_{K,r}(t)^{-1}=g_{K,r}(-t)=g_{K,r^{-1}}(t)\).
\end{theorem}
\begin{proof}
Todas las matrices son diagonales. La entrada \(i\) de la primera
igualdad es \((ab)^{tv_i}=a^{tv_i}b^{tv_i}\), puesto que
\(\log(ab)=\log a+\log b\) para bases positivas. Cada entrada
es estrictamente positiva, por lo que la matriz es invertible, y
la sustitución \(b\mapsto b^{-1}\) prueba la segunda igualdad.
La tercera se obtiene de \(r^{(s+t)v_i}=r^{sv_i}r^{tv_i}\).
El determinante es el producto de las doce entradas:
\[
 \prod_{i=1}^{12}r^{tv_i}
 =r^{t\sum_i v_i}=r^0=1.
\]
La fórmula de la inversa es nuevamente la inversión de cada entrada.
\end{proof}

\begin{corollary}[Composición de ambivalencia y espejo regional]
\label{hrr:composicion}
Para todo \(n\geq2\),
\begin{equation}
 g_{K,\Theta_{\pi,n}}(t)
 g_{K,\Theta_{e,n}}(t)^{-1}
 =g_{K,\pi/e}(t).
 \label{hrr:composicion-finita}
\end{equation}
En el límite estable,
\begin{equation}
 g_{K,\pi/\varphi^2}(t)
 =g_{K,\pi/e}(t)\,g_{K,e/\varphi^2}(t).
 \label{hrr:composicion-limite}
\end{equation}
\end{corollary}
\begin{proof}
Las dos marcas finitas son positivas para \(n\geq2\). Aplicando
\eqref{hrr:cociente-flujo} y \eqref{hrr:cociente-exacto}, su cociente
de matrices es el flujo de la marca \(\pi/e\). Para la segunda
igualdad basta usar
\((\pi/e)(e/\varphi^2)=\pi/\varphi^2\) en
\eqref{hrr:producto}. También se obtiene como límite de la identidad
finita, porque la exponencial y el logaritmo son continuos en el
dominio de bases positivas.
\end{proof}

La igualdad compara deformaciones con una misma dirección \(u_K\).
El intercambio especular interno de los canales regionales conserva
su propio dominio sobre el estado enriquecido; la razón \(\pi/e\)
es su lectura escalar posterior. El operador de holonomía \(\Gamma_9\)
continúa actuando sobre fase y memoria. Estos dominios determinan el
sentido de la composición y preservan el registro que acompaña al retorno.

\subsection{Transporte de menores e invariantes exteriores}

Sea \(X\) una matriz real de rango \(k\) con doce columnas y
\(\Delta_I(X)\) el menor de columnas indicado por un conjunto
\(I\) de \(k\) índices. La acción por la derecha de
\(g_{K,r}(t)\) escala la columna \(i\) por \(r^{tv_i}\).

\begin{proposition}[Signos, estratos e invariantes de cocientes]
\label{hrr:plucker}
Se cumple
\begin{equation}
 \Delta_I\bigl(Xg_{K,r}(t)\bigr)
 =r^{t\sum_{i\in I}v_i}\Delta_I(X).
 \label{hrr:plucker-formula}
\end{equation}
Se conservan los signos y el patrón de anulación de todos los menores.
Si \(I,J,P,Q\) satisfacen
\(\mathbf1_I+\mathbf1_J=\mathbf1_P+\mathbf1_Q\), entonces,
donde \(\Delta_P(X)\Delta_Q(X)\neq0\), el cociente
\begin{equation}
 \frac{\Delta_I\Delta_J}{\Delta_P\Delta_Q}
 \label{hrr:cociente-menores}
\end{equation}
es invariante bajo el flujo.
\end{proposition}
\begin{proof}
La multilinealidad del determinante multiplica el menor por el producto
de los factores de sus columnas, lo que prueba
\eqref{hrr:plucker-formula}. Ese producto es positivo y, por tanto,
conserva el signo y la nulidad del menor. En
\eqref{hrr:cociente-menores}, el exponente total es
\[
 t\sum_{i=1}^{12}
 (\mathbf1_I(i)+\mathbf1_J(i)-\mathbf1_P(i)-\mathbf1_Q(i))v_i=0.
\]
El denominador permanece distinto de cero y la cancelación es exacta.
\end{proof}

Para \(k=2\), una aplicación concreta es
\[
 \frac{\Delta_{12}\Delta_{34}}{\Delta_{13}\Delta_{24}}.
\]
Cada uno de los cuatro índices aparece una vez en el numerador y una
vez en el denominador. En una configuración donde los menores del
denominador sean distintos de cero, esta cantidad permanece fija,
aunque cada menor individual cambie. La conservación exterior posee
así observables explícitos, adicionales al determinante global.

\subsection{Elevación reticular y dualidad}

En la realización reticular marcada del capítulo precedente se utiliza
la suma ortogonal de doce planos. Escribamos
\[
 G_{K,r}(t)=\bigoplus_{i=1}^{12}r^{tv_i}I_2.
\]
Si \(\Lambda\) es la red euclídea autodual de referencia ya obtenida,
definimos \(\Lambda_{r,t}=G_{K,r}(t)\Lambda\).

\begin{proposition}[Covolumen y reciprocidad de las redes deformadas]
\label{hrr:redes}
El covolumen de \(\Lambda_{r,t}\) coincide con el de \(\Lambda\), y
\begin{equation}
 \Lambda_{r,t}^{*}=\Lambda_{r,-t}.
 \label{hrr:red-dual}
\end{equation}
La factorización \eqref{hrr:composicion-limite} se eleva a estos
operadores de veinticuatro dimensiones.
\end{proposition}
\begin{proof}
Cada radio aparece dos veces en el determinante. Por tanto,
\[
 \det G_{K,r}(t)=\prod_i r^{2tv_i}=1.
\]
El cambio lineal multiplica el covolumen por el valor absoluto del
determinante, por lo que lo conserva. Para cualquier operador invertible
\(G\), la condición \(y\in(G\Lambda)^*\) equivale a
\(\langle y,Gx\rangle\in\mathbb Z\) para todo \(x\in\Lambda\).
Equivale a \(G^{\mathsf T}y\in\Lambda^*\), y por ello
\((G\Lambda)^*=G^{-\mathsf T}\Lambda^*\). Aquí
\(\Lambda^*=\Lambda\), \(G\) es diagonal positivo y
\(G^{-\mathsf T}=G_{K,r}(-t)\), lo que demuestra
\eqref{hrr:red-dual}. La factorización se verifica sobre cada bloque
de dos dimensiones exactamente como en \eqref{hrr:producto}.
\end{proof}

Esta conclusión conserva el covolumen y relaciona la deformación con
su dual. La integralidad de los productos escalares y la paridad de
las normas son propiedades adicionales de una red. La reciprocidad
espectral theta--Mellin utiliza precisamente la dualidad demostrada,
junto con la sumabilidad gaussiana y la fórmula de Poisson desarrolladas
en su capítulo propio.

\section{Diagonal normalizada y restitución de la orientación radial}
\label{hrr:diagonal}

\subsection{Separación de escala y forma}

La carta radial antecedente parte de las coordenadas angulares generadas
\(A>0\) y \(C^*\), en el dominio \(|C^*|<A\). Definimos
\[
 \xi=\frac{C^*}{A},\qquad
 \eta=\operatorname{artanh}\xi.
\]
La variable \(\eta\) designa la anisotropía de la forma. Su significado
es distinto del de la mantisa de retorno \(\eta_{\mathrm{ret}}\)
que interviene en la sección de acción.

Sea \(s>0\) la escala areal producida por esa sección. En la realización
de la elipse de acción, \(s=2\hbar\). Los semiejes positivos son
\begin{equation}
 a=\sqrt{s}\,e^{\eta/2},\qquad
 b=\sqrt{s}\,e^{-\eta/2},\qquad ab=s.
 \label{hrr:semiejes}
\end{equation}
La reciprocidad intercambia \(a\) y \(b\). Su coordenada radial es
\(R_*=\sqrt{b/a}=e^{-\eta/2}\). La longitud considerada a
continuación es la diagonal del rectángulo de lados \(a,b\):
\[
 d=\sqrt{a^2+b^2},\qquad D=\frac{d}{\sqrt{ab}}.
\]
La normalización elimina la escala areal y conserva la forma del
rectángulo hasta el intercambio de sus lados.

\begin{theorem}[Diagonal, anisotropía y reciprocidad]
\label{hrr:diagonal-teorema}
En el dominio anterior,
\begin{equation}
 D^2=2\cosh\eta=R_*^2+R_*^{-2},
 \qquad
 D=\frac{\sqrt2}{(1-\xi^2)^{1/4}}.
 \label{hrr:diagonal-formula}
\end{equation}
Se cumple \(D\geq\sqrt2\), y son equivalentes
\[
 D=\sqrt2,\qquad \eta=0,\qquad \xi=0,\qquad R_*=1,
 \qquad a=b.
\]
Para \(D>\sqrt2\), el dato adicional
\(\varepsilon=\operatorname{sgn}\eta\in\{-1,+1\}\) determina
de manera única
\begin{align}
 \eta&=\varepsilon\operatorname{arcosh}(D^2/2),
 \label{hrr:eta-recuperada}\\
 \xi&=\varepsilon\sqrt{1-4/D^4},
 \label{hrr:xi-recuperada}\\
 R_*&=\exp\!\left[-\frac{\varepsilon}{2}
                   \operatorname{arcosh}(D^2/2)\right].
 \label{hrr:radio-recuperado}
\end{align}
\end{theorem}
\begin{proof}
Dividiendo \(a^2+b^2\) por \(ab=s\) en
\eqref{hrr:semiejes}, resulta
\[
 D^2=e^{\eta}+e^{-\eta}=2\cosh\eta.
\]
Como \(R_*^2=e^{-\eta}\), la suma es también
\(R_*^2+R_*^{-2}\). De \(\xi=\tanh\eta\) y
\(1-\tanh^2\eta=\cosh^{-2}\eta\), con
\(\cosh\eta>0\), se obtiene
\(\cosh\eta=(1-\xi^2)^{-1/2}\). La raíz positiva proporciona
la segunda fórmula de \eqref{hrr:diagonal-formula}.

La desigualdad \(\cosh\eta\geq1\), con igualdad sólo para
\(\eta=0\), prueba la cota y sus equivalencias mediante las
definiciones. Para \(D>\sqrt2\), \(\cosh\) determina
\(|\eta|=\operatorname{arcosh}(D^2/2)\), y la orientación
\(\varepsilon\) determina su signo. Además,
\[
 \xi^2=1-\cosh^{-2}\eta=1-\frac4{D^4}.
\]
La monotonía de \(\tanh\) implica
\(\operatorname{sgn}\xi=\operatorname{sgn}\eta\), de donde
\eqref{hrr:xi-recuperada}. La sustitución de
\eqref{hrr:eta-recuperada} en \(R_*=e^{-\eta/2}\) prueba la
última fórmula. Cada paso es unívoco en el dominio indicado.
\end{proof}

\begin{corollary}[Recuperación exacta de los semiejes]
\label{hrr:semiejes-recuperacion}
Los datos \((s,D,\varepsilon)\), con \(s>0\) y
\(D>\sqrt2\), recuperan los semiejes ordenados por
\begin{equation}
 a^2=\frac{s}{2}\left(D^2+\varepsilon\sqrt{D^4-4}\right),
 \qquad
 b^2=\frac{s}{2}\left(D^2-\varepsilon\sqrt{D^4-4}\right).
 \label{hrr:semiejes-radicales}
\end{equation}
Si \(D=\sqrt2\), la recuperación es \(a=b=\sqrt{s}\).
\end{corollary}
\begin{proof}
Las ecuaciones \(a^2+b^2=sD^2\) y \(a^2b^2=s^2\)
determinan \(a^2,b^2\) como las raíces positivas de
\(z^2-sD^2z+s^2=0\). Sus raíces son las dos cantidades
\(s(D^2\pm\sqrt{D^4-4})/2\). El signo de
\(a-b\) coincide con el de \(\eta\), de modo que
\(\varepsilon\) asigna cada raíz a su semieje. Las raíces son
positivas porque \(D^2>\sqrt{D^4-4}\). El caso de igualdad
produce una raíz doble \(s\).
\end{proof}

\subsection{Ejemplo exacto de dos realizaciones recíprocas}

Tomemos una forma de prueba con \(\xi=3/5\). Este ejemplo verifica
el lector geométrico; la asignación concreta de una sección HMT se
obtiene de sus propias coordenadas generadas. Se calcula
\[
 \eta=\operatorname{artanh}(3/5)=\log2,\qquad
 R_*=\frac1{\sqrt2},\qquad
 D^2=2+\frac12=\frac52.
\]
Para escala \(s\), los semiejes son
\(a=\sqrt{2s}\) y \(b=\sqrt{s/2}\). La forma recíproca tiene
\(\xi=-3/5\), \(\eta=-\log2\), \(R_*=\sqrt2\) y los
semiejes intercambiados. Ambas poseen la misma diagonal normalizada
\(D=\sqrt{5/2}\) y el mismo producto \(ab=s\). El dato de
orientación distingue exactamente esas dos realizaciones.

La diagonal representa por tanto una intensidad de anisotropía. El
registro orientado permite reconstruir la forma ordenada. La escala
areal, por su parte, recupera su tamaño. Esta separación especifica
qué dato conserva cada lector y qué dato aporta el archivo.

\subsection{Cartas de Poincaré y Klein}

La carta de Poincaré del diámetro radial y su carta de Klein son
\begin{equation}
 w=\frac{R_*^2-1}{R_*^2+1}
   =-\tanh(\eta/2),\qquad
 x_{\mathrm K}=\frac{2w}{1+w^2}=-\tanh\eta=-\xi.
 \label{hrr:cartas}
\end{equation}
Los dominios \(R_*>0\), \(-1<w<1\) y \(-1<x_{\mathrm K}<1\)
se corresponden mediante estas aplicaciones.

\begin{proposition}[Reciprocidad radial en las dos cartas]
\label{hrr:reciprocidad-cartas}
La inversión \(R_*\mapsto R_*^{-1}\) actúa por
\[
 \eta\mapsto-\eta,\qquad w\mapsto-w,
 \qquad x_{\mathrm K}\mapsto-x_{\mathrm K}.
\]
La diagonal y la escala areal permanecen invariantes. El par
\((D,\varepsilon)\) recupera ambas coordenadas de disco mediante
\eqref{hrr:eta-recuperada} y \eqref{hrr:cartas}.
\end{proposition}
\begin{proof}
La identidad \(R_*=e^{-\eta/2}\) transforma la inversión en el
cambio de signo de \(\eta\). Las dos funciones hiperbólicas de
\eqref{hrr:cartas} son impares, por lo que cambian de signo.
La fórmula \(D^2=R_*^2+R_*^{-2}\) es invariante y el producto
de los semiejes permanece \(s\). Las fórmulas inversas del
teorema~\ref{hrr:diagonal-teorema} completan la recuperación.
\end{proof}

En el ejemplo \(\xi=3/5\), se obtiene
\(w=-1/3\) y \(x_{\mathrm K}=-3/5\). La forma recíproca
tiene \(w=1/3\) y \(x_{\mathrm K}=3/5\). La coincidencia de
la diagonal acompaña así a dos posiciones opuestas del mismo diámetro.

\subsection{Inversión radial y reflexión en la carta compleja}

El radio positivo también interviene en la carta compleja de
reciprocidad. Para \(r>0\) y un argumento \(\theta\), sea
\[
 z=re^{i\theta},\qquad s(r,\theta)=\frac1{1-z},
 \qquad z\neq1.
\]
El dominio excluye únicamente el polo de esta carta. La inversión
radial conserva \(\theta\) y sustituye \(r\) por \(r^{-1}\).

\begin{proposition}[Reflexión respecto de la recta de parte real media]
\label{hrr:reflexion-compleja}
En el dominio anterior,
\begin{align}
 s(r^{-1},\theta)&=1-\overline{s(r,\theta)},
 \label{hrr:reflexion}\\
 \operatorname{Re}s(r,\theta)-\frac12
 &=\frac{1-r^2}{2|1-re^{i\theta}|^2}.
 \label{hrr:reflexion-parte-real}
\end{align}
Por tanto, el círculo \(r=1\), salvo el polo, se transforma en la
recta \(\operatorname{Re}s=1/2\). La inversión radial se realiza
como reflexión respecto de esa recta vertical.
\end{proposition}
\begin{proof}
El conjugado de \(s(r,\theta)\) es
\((1-re^{-i\theta})^{-1}\). Por cálculo directo,
\[
 1-\frac1{1-re^{-i\theta}}
 =\frac{-re^{-i\theta}}{1-re^{-i\theta}}
 =\frac1{1-r^{-1}e^{i\theta}},
\]
lo que demuestra \eqref{hrr:reflexion}. Para la segunda identidad,
se multiplica numerador y denominador de \(s\) por
\(1-re^{-i\theta}\). Su parte real es
\[
 \frac{1-r\cos\theta}{1-2r\cos\theta+r^2}.
\]
Al restar \(1/2\), el numerador se reduce a \((1-r^2)/2\).
El denominador es positivo en el dominio. De aquí se deduce que
\(\operatorname{Re}s=1/2\) equivale a \(r=1\). Finalmente,
si \(s=x+iy\), entonces \(1-\overline s=(1-x)+iy\), que
es exactamente la reflexión respecto de la recta vertical indicada.
\end{proof}

Para \(\theta=\pi\), los radios \(r=2\) y \(r=1/2\)
producen respectivamente \(s=1/3\) y \(s=2/3\). El punto
\(r=1\) produce \(s=1/2\). Este ejemplo hace visible la
reciprocidad mediante tres valores racionales de la carta.

La transformación demostrada pertenece a la geometría de la carta
compleja. Su aplicación a un operador espectral exige conservar el
operador y su dominio, además de esa carta. Del mismo modo, la carta
radial \(R_*\) conserva la sección angular que la genera. La igualdad
de una operación de inversión en dos realizaciones proporciona la
correspondencia escrita; la identificación de sus objetos antecedentes
se mantiene mediante los transportes ya construidos.

\section{Recuperación del flujo y compatibilidad de los lectores}
\label{hrr:archivo}

La aplicación de estos lectores a una dirección de flujo exige conservar
la dirección generada, además de sus valores escalares. El archivo
reducido construido previamente posee operadores \(M\) y \(L\)
con \(LM=P_{11}\). Si \(m=MK\), entonces
\begin{equation}
 u_K=P_3Lm.
 \label{hrr:archivo-direccion}
\end{equation}
Esta igualdad mantiene la acción de cada operador y permite componer
la recuperación con el flujo.

\begin{proposition}[Reconstrucción del flujo desde el archivo reducido]
\label{hrr:archivo-flujo}
Con los operadores y la marcación precedentes, el archivo \(m\) y
una marca positiva \(r\) determinan
\[
 g_{K,r}(t)=
 \operatorname{diag}\left(
 r^{\,t(P_3Lm)_1/\|P_3Lm\|},\ldots,
 r^{\,t(P_3Lm)_{12}/\|P_3Lm\|}\right).
\]
Se conservan todas las identidades de composición de
\eqref{hrr:producto}--\eqref{hrr:determinante}.
\end{proposition}
\begin{proof}
De \(m=MK\) y \(LM=P_{11}\) se deduce
\(P_3Lm=P_3LMK=P_3P_{11}K=u_K\).
La dirección es distinta de cero por la propiedad ya demostrada del
registro generado. La normalización recupera exactamente \(v\),
y su sustitución en \eqref{hrr:g-definicion} recupera cada entrada
del flujo. Las identidades dependen únicamente de esa dirección y
de la positividad de las marcas, por lo que permanecen válidas.
\end{proof}

La identidad de archivo recupera la dirección de la deformación.
Los lectores regionales aportan sus marcas. El lector diagonal de
la forma radial conserva su anisotropía absoluta, y el registro
orientado recupera el signo. Se obtiene una composición en la cual
la invariancia y la reversibilidad se asignan a datos explícitos:
\[
 \begin{gathered}
 \text{retornos de hoja}
 \longrightarrow(\Theta_{\pi,n},\Theta_{e,n})
 \longrightarrow\pi/e,\\
 \text{archivo de }K\longrightarrow u_K
 \longrightarrow g_{K,r}(t)
 \longrightarrow\text{covolumen y menores transportados},\\
 (s,D,\varepsilon)\longrightarrow(a,b)
 \longleftrightarrow(R_*,w,x_{\mathrm K}).
 \end{gathered}
\]
La primera línea distingue el soporte común de dos marcas. La segunda
expresa la función operatoria del registro. La tercera especifica la
información geométrica necesaria para la inversión. El estado enriquecido
mantiene la fase, las vueltas, los acarreos y las incidencias que permiten
componer estos lectores dentro de su genealogía.

\section{Realización temporal y dominio de interpretación física}

La holonomía nonádica conserva su acción sobre fase y memoria; el retorno
de fase incrementa el registro de vueltas. Los recuentos del envolvente
modal y las razones de marcas se evalúan sobre realizaciones posteriores.
La reciprocidad radial actúa sobre una carta de forma y su orientación.
La articulación demostrada reúne estas operaciones mediante lectores
compatibles, conservando la distinción entre sus parámetros.

La relación \(\pi/\varphi^2\) comparece como razón marcada de retornos
y como escala del flujo conservativo. La razón \(\pi/e\) comparece
como cociente de marcas sobre el mismo soporte. Ambas satisfacen la
factorización exacta \eqref{hrr:composicion-limite}. La intervención
de estas razones en una realización gravitatoria o cosmológica conserva
las condiciones del correspondiente lector físico: identificación del
reloj, magnitudes dimensionales, respuesta de frontera y dominio geométrico.
El resultado de este capítulo pertenece a la composición matemática de
los lectores generados. Sus aplicaciones físicas utilizan después la
realización que les corresponde.
